% !PRINTED 372
\weberpage{372}
Bilden wir die beiden Formen
\[
\begin{aligned}
\lambda &= a_1 a_2\,t_1 + \frac{b-\sqrt D}{2}\,t_2, \\
\lambda' &= a\,t_1 + \frac{b-\sqrt D}{2}\,t_2,
\end{aligned}
\]
so folgt:
\[
\begin{aligned}
N(\lambda) &= a_1 a_2 (a_1 a_2\,t_1^2 + b\,t_1 t_2 + c\,t_2^2), \\
N(\lambda') &= a (a\,t_1^2 + b\,t_1 t_2 + a'c\,t_2^2).
\end{aligned}
\]
Die beiden Formen $(a_1 a_2,\ b,\ c)$, $(a,\ b,\ a'c)$ sind primitiv (weil ein gemeinsamer Primteiler von $a_1 a_2$, $b$, $c$ wegen (15) in $Q$ aufgehen müßte), und es folgt für die absolute Norm:
\[
\begin{aligned}
N_a(\lambda) &= a_1 a_2 = N(\mathfrak{a}\mathfrak{b}), \\
N_a(\lambda') &= a;
\end{aligned}
\]
da aber $\lambda'$ durch $\mathfrak{a}\mathfrak{b}$ teilbar ist, so muß $a$ durch $a_1 a_2$ teilbar und mithin $a'=1$ sein.

Hieraus ergibt sich:

3. \gesperrt{Sind $a_1$, $a_2$, $b$ ohne gemeinsamen Teiler und}
\[
D = b^2 - 4\,a_1 a_2 c,
\]
\gesperrt{so ist die Form}
\[
T = (a_1 a_2,\ b,\ c)
\]
\gesperrt{aus den beiden Formen}
\[
U = (a_1,\ b,\ c\,a_2), \quad V = (a_2,\ b,\ c\,a_1)
\]
\gesperrt{komponiert.}

Wendet man dies auf zwei entgegengesetzte Formen
\[
U = (a,\ b,\ c), \quad V = (c,\ b,\ a)
\]
an, so ergibt sich
\[
T = (ac,\ b,\ 1),
\]
und dies ist mit der Hauptform $\left(1,\ 0,\ \dfrac{-D}{4}\right)$ oder $\left(1,\ 1,\ \dfrac{1-D}{4}\right)$ äquivalent.

4. \gesperrt{Zwei entgegengesetzte Klassen geben komponiert die Hauptklasse.}

Um eine Klasse $A$ wiederholt mit sich selbst zusammenzusetzen, wähle man $\mathfrak{a}$ in $A$ so, daß es relativ prim zu $D$ ist, und bezeichne mit $(a,\ b,\ c)$ die entsprechende Form. Dann ist $a$ relativ prim zu $b$. Denn wäre $d$ ein Teiler von $a$ und $b$, so wäre $d$ auch in $D$ enthalten, also relativ prim zu $\mathfrak{a}$. Es wäre also auch $a:d$ durch $a$ teilbar, was, wenn $d>1$ ist, der Definition von $a$ widerspricht.
